\chapter{Evidence for Model Adequacy}
\index{evidence!model adequacy|textbf}

\label{ch:evidence-for-model-adequacy}

\chapterstatus{evidence framework and retained-status interpretation; no new
scientific validation or uncertainty-quantification result is claimed.}

Chapter~\ref{ch:training-evidence-vocabulary} defined bounded claims,
oracles, metrics, tolerances, dispositions, and the five evidence classes used
throughout the monograph. This chapter applies that vocabulary to the reduction
program's verification inventory, convergence questions, validation targets,
uncertainty sources, provenance, and negative results.

\section{Verification in the reduction program}

\subsection{Software verification}
\index{verification!software}


Software-verification evidence uses public contracts, schemas, exact language
semantics, or otherwise independently stated behavior as its oracle. For the
present program, relevant subjects include

\begin{itemize}
  \item construction of represented-operator records and their intrinsic
    invariants;
  \item basis, geometry, energy-reference, and spin metadata;
  \item Hermiticity and symmetry analyzers;
  \item compatibility checks that precede arithmetic;
  \item signed difference and residual conventions;
  \item serializer field vocabularies and exact promised round trips;
  \item tight-binding limiting cases and Fourier conventions;
  \item alignment recovery under known synthetic unitary transformations;
  \item continuum-embedding behavior for synthetic envelopes; and
  \item external-tool lifecycle, stream capture, and provenance manifests.
\end{itemize}

Tests should use the supported public surface and an oracle that does not simply
reproduce the production algorithm. Exact contracts use exact acceptance;
numerical contracts use explicit tolerances justified by the operation and
scale. Helpers and fixtures remain distinguishable from evidence-bearing cases.

A round trip establishes only the promised wire representation. A compatibility
result establishes only that the represented prerequisites checked by that
operation agree. Neither proves that recorded provenance is truthful or that the
physical model is adequate.

\subsection{Numerical verification}
\index{verification!numerical}


Numerical verification asks whether the implementation approximates the stated
mathematics. Controlled cases are especially important before first-principles
artifacts are available. Examples include

\begin{itemize}
  \item exact recovery of an operator after a known unitary gauge change;
  \item principal angles for analytically constructed subspaces;
  \item residual norms derived independently for small matrices;
  \item Hermitian tight-binding operators in zero-hopping or symmetry-limited
    cases;
  \item Fourier transforms checked against finite hand-constructed lattices;
  \item effective-mass stencils applied to an analytic quadratic dispersion;
  \item continuum solvers checked against solvable potentials where the model
    assumptions match; and
  \item embedding maps checked with synthetic normalized envelopes.
\end{itemize}

Synthetic data are labelled as such. Exact recovery on a controlled example is
strong evidence about the corresponding mathematical implementation, but it is
not a first-principles calculation or physical validation.

\subsection{Compatibility witnesses and separation certificates}
\index{evidence!common witness}
\index{evidence!separation certificate}

Chapter~\ref{ch:admissible-sets-certificates} defines compatibility as the
intersection of thresholded parameter sets. The evidence supporting an
intersection differs logically from the evidence supporting separation.

A retained parameter is a \emph{common witness} only when independently
evaluated losses show that the same parameter satisfies every declared
criterion. It proves existence within the frozen family and thresholds. It does
not prove uniqueness, stability, transferability, or scientific validation.

A \emph{separation certificate} requires a lower bound that applies to the
complete admissible sets, including every applicable alignment branch and
parameter-domain boundary. Two numerically optimized points supply only an
upper bound on the distance between sets. A failed grid or local search supplies
neither a common witness nor a positive global lower bound. When neither form
of evidence is available, the correct disposition is unresolved.

The certificate's premises are part of its evidence. For a quadratic
axis-reduction argument these include the retained quadratic coefficients,
positive curvature, symmetry, zero or otherwise accounted-for cross terms,
the location of the centers, and the absence or explicit treatment of domain
clipping. A verifier that recomputes a final scalar without authenticating
these premises has not independently reconstructed the full certificate.

Threshold sensitivity has a separate role. A prospectively frozen threshold
defines the original decision; a later sweep may reveal how close that decision
lies to a transition. The latter is post-hoc sensitivity evidence unless it was
itself frozen in advance. It must not be presented as physical calibration of
the threshold or as uncertainty quantification.

\subsection{Completed controlled benchmark inventory}
\index{evidence!controlled synthetic benchmarks}


Appendix~\ref{app:one-dimensional-reduction} supplies a calculated example of
this evidence class. It compares independently refined plane-wave and
finite-difference parents, checks selected energies against Mathieu
characteristic values, and compares direct and Wannier90 localization using the
same spectral and overlap data. Discretization, subspace conditioning, gauge,
localization, serialization, and hopping truncation remain separate error
sources. Nonconverged localization attempts are retained as diagnostic evidence
rather than relabelled as successful calculations.

Section~\ref{sec:appendix-g-m3} separately records the canonical M3
constrained-admissible-set extension. M3 supplies one retained common witness and one analytic positive
separation certificate for two frozen threshold pairs in a bounded synthetic
rank-two family. Its public and standalone verifiers reconstruct the role,
loss, threshold, finite-alignment, and certificate contracts. The associated
threshold-sensitivity package is post-hoc exploratory evidence and does not
calibrate a physical acceptance threshold.

Appendix~\ref{app:two-dimensional-reduction} extends the controlled
periodic benchmark to two dimensions. It verifies separable and coupled scalar
parents, composite-band projection with the corrected polar gauge, mass-tensor
recovery, real-space hopping reconstruction, common-estimator localization
comparisons, and distinct QWZ, Hofstadter, and Haldane topology controls.
Mesh, cutoff, embedding, hopping truncation, serialization, topology,
optimizer-basin, and continuation effects remain separate evidence axes.

The bounded localization study is complete as negative numerical evidence.
Across the 256 frozen start/interface trajectories, the final effective
endpoints comprise 196 native convergences and 60 iteration-limit
nonconvergences. The frozen mesh-convergence and repeated-basin criteria are
not supported. The censored convergence-time regression is exploratory and
descriptive. Only one frozen trial-projection family was executed, so
projection-family sensitivity is outside the scope of the closed benchmark.
These results establish neither global optimization nor general Wannier
convergence.

Appendix~\ref{app:impurity-model-classes} separately verifies spin-space
contracts and one-dimensional pristine--defect extraction, including blind
alignment, an independent extraction route, a finite-rank spectral oracle, and
continuum refinement. Those accepted synthetic records do not establish a
two-dimensional defect result or material validation. The proposed
two-dimensional extension is therefore kept outside the completed evidence
inventory in Appendix~\ref{app:two-dimensional-defect-extraction}.

\subsection{Finite-setting convergence}
\index{convergence!finite-setting}

Chapter~\ref{ch:training-evidence-vocabulary} explains why stability over a
finite tested sequence does not by itself bound an infinite-setting limit. In
the present program, cutoff and mesh errors may additionally interact. For an
accepted candidate pair
$(E_{\ast},K_{\ast})$ and higher guards $(E_{+},K_{+})$, a finite two-factor
contrast is
\begin{equation}
  I_q=
  \left|
    q(E_{+},K_{+})-q(E_{+},K_{\ast})
    -q(E_{\ast},K_{+})+q(E_{\ast},K_{\ast})
  \right|.
\end{equation}
A small value detects limited interaction in that tested rectangle; it does not
prove global separability of cutoff and sampling errors.

\subsection{Derived quantities and error amplification}
\index{derived quantity}
\index{error!amplification}


A pointwise energy tolerance need not control a derivative. For a central
second-difference estimate,
\begin{equation}
  E''(k_0)
  \approx
  \frac{E(k_0+h)-2E(k_0)+E(k_0-h)}{h^2},
\end{equation}
pointwise energy perturbations are amplified by a scale proportional to
$h^{-2}$. Effective-mass verification must therefore vary the stencil or fit
window, report conditioning, identify the selected valley and axes, and check
stability under at least one smaller scale. A converged band plot is not an
effective-mass convergence study.

The same principle applies to subtraction of two large aligned operators,
normalization by a small reference scale, and composition of successive
reduction maps. Numerical cancellation, conditioning, and nonfinite
intermediates must be evaluated at the level of the derived quantity.

\section{Validation and uncertainty}

\subsection{Scientific validation}
\index{validation!scientific}


Scientific validation compares a declared model output with independent
physical evidence appropriate to the intended use. Potential bulk-silicon
references include experimental structural or band-edge measurements and
higher-level theoretical results, but each comparison must match temperature,
geometry, relativistic treatment, observable definition, and uncertainty as far
as the claim requires.

The PBE Kohn--Sham parent is not fitted silently to the experimental silicon gap.
Its disagreement with experiment belongs to parent-model assessment. A
Wannier or tight-binding model is first evaluated against that selected parent
to measure reduction fidelity. If the reduced model is then compared with
experiment, the account must not attribute the parent's discrepancy to the
reduction or claim cancellation as validation without a declared rationale.

\citationtodo{high priority}{Check the silicon-specific PBE gap statement with
a settings-relevant PBE benchmark and an experimental or evaluated silicon
reference. Candidate sources: Paier et al. (2006), provisional key
\texttt{paier2006}, and Yu--Cardona (2010),
\texttt{yuCardona2010}. Perdew--Levy and Sham--Schluter, keys
\texttt{perdewLevy1983} and \texttt{shamSchluter1983}, address why the exact
Kohn--Sham and physical gaps require care.}

For dopants, validation references must match the charge-state interpretation,
periodic or isolated setting, spin and SOC branch, band-edge reference, and
observable definition. A neutral periodic P:Si calculation cannot be validated
as though it were directly the potential of an isolated ionized donor.

\citationtodo{high priority}{Check this charge-state and boundary-condition
matching rule against defect-supercell literature. Candidate: Freysoldt et al.
(2014), bibliography key \texttt{freysoldt2014} (already present).}

\subsection{Uncertainty quantification}
\index{uncertainty quantification}


Relevant uncertainty sources can include

\begin{itemize}
  \item parent-model choice and external-reference uncertainty;
  \item cutoff, mesh, cell-size, and solver discretization;
  \item lattice-fit and valley-location uncertainty;
  \item effective-mass fit-window and stencil sensitivity;
  \item Wannier windows, projections, gauge, and localization choices;
  \item alignment and scalar energy-reference uncertainty;
  \item parameter non-identifiability and optimization variability;
  \item model-class and real-space truncation sensitivity; and
  \item continuum discretization, boundary, and embedding uncertainty.
\end{itemize}

These sources are not combined automatically. A valid propagation statement
requires compatible quantities, an explicit dependence model, and treatment of
correlation where material. Until then, the sources are reported separately as
sensitivity or bounded numerical evidence.

\section{Model development and provenance}

\subsection{Training, validation, and leakage}
\index{training data}
\index{validation data}
\index{data leakage}


Model fitting uses only the frozen training set. Withheld points, band-edge
quantities, operator blocks, or dopant observables designated for validation
must not influence parameter selection, model-class expansion, stopping rules,
or tolerance adjustment. If a validation failure motivates a new model class,
that class and a new untouched validation design must be recorded explicitly;
the original result remains part of the evidence history.

\citationtodo{medium priority}{Check the validation-leakage rule and the need for
a new untouched validation design after model redesign against model-selection
literature. Candidate: Cawley and Talbot (2010), provisional key
\texttt{cawley2010}; the candidate record is present for this editorial note.}

Visual agreement can support diagnosis but is not an acceptance rule. Plots
must expose the data identity, axes, units, reference, and residual definition,
while machine-readable results retain the exact values used for the decision.

\subsection{Provenance and reproducibility}
\index{provenance}
\index{reproducibility}


A result is reproducible only relative to retained identities and declared
conditions. The minimum record depends on the operation but generally includes

\begin{itemize}
  \item immutable input and source identities;
  \item code revision and dependency or executable versions;
  \item physical and numerical settings;
  \item exact commands and environment information;
  \item artifact manifests and checksums;
  \item metric and tolerance definitions;
  \item structured outcomes and warnings; and
  \item locations of externally retained large artifacts.
\end{itemize}

Reproduction does not promote the evidence class. Reproducing a tutorial
calculation reproduces tutorial behavior; reproducing a synthetic test
reproduces software or numerical verification; reproducing a production
calculation reproduces that calculation under its recorded conditions.

\citationtodo{medium priority}{Check the minimum reproducibility record against
established computational-research guidance. Candidate: Sandve et al. (2013),
provisional key \texttt{sandve2013}; the candidate record is present for this
editorial note.}

\subsection{Current execution-status example}

The repository retains a direct bulk-silicon bootstrap campaign in which nine
SCF and nine linked NSCF Quantum ESPRESSO invocations exited successfully and
reported completion. The owning disposition classifies these observations as
bootstrap scientific-harness development evidence. They are not a canonical
scientific workflow record and do not accept a final cutoff, mesh, lattice
parameter, gap, mass, infinite-basis result, or scientific validation claim.

This example demonstrates why process completion and evidentiary acceptance are
separate. The observations may inform integration and future workflow design
without being discarded or promoted into a production result.

\section{Failure and claim discipline}

\subsection{Failure and negative evidence}
\index{negative evidence}
\index{failure}


A failed acceptance criterion is retained rather than weakened. The diagnosis
may identify an implementation defect, inadequate numerical resolution,
incompatible representation, insufficient model-class expressiveness, or
mismatch between a model and its intended physical use. These dispositions have
different remedies.

Similarly, a well-executed negative result can be scientifically useful: no
model in a frozen hierarchy may satisfy both spectral and operator criteria; no
finite tested crossover radius may satisfy continuum criteria; two reduction
paths may fail to commute within tolerance. The evidence must report the tested
domain and must not generalize beyond it.

\subsection{Claim-to-evidence discipline}
\index{claim!evidence discipline}


Later results are included only when each statement can point to an applicable
retained record. Software verification supports contract claims. Numerical
verification supports stated mathematical or convergence claims. Scientific
validation supports adequacy for a declared physical use. UQ supports only the
uncertainties it actually identifies and propagates.

This discipline is not administrative overhead added after computation. It is
how the project prevents a passing test, a successful executable, a smooth plot,
or an aligned matrix from being mistaken for a stronger scientific conclusion.
